% Chapter introduction: why the topic arose, who first asked, what it answers.
% STATUS: rewritten 2026-10-08 (issue #678); every claim is cited in its section.

Chapter~\ref{ch:origins} followed the product rule from Leibniz to Hochschild.
By 1945 the rule was a definition on abstract algebras. This chapter asks what
one can do with that definition. Three groups of mathematicians gave answers,
and each group studied a different object.

The first group studied differential equations. Ritt wanted a field that comes
with a derivative, so that the algebra of such fields could be built like the
algebra of ordinary fields \cite[p.~1]{Ritt1950}. His main object was the set
of elements with zero derivative, called the constants.
Section~\ref{sec:basic-theory:constants-kernels} returns to them.

The second group studied Lie algebras. Jacobson saw that the derivations of an
algebra are closed under the commutator, so they form a Lie algebra
\cite[p.~207]{Jacobson1937}. Hochschild then showed what this Lie algebra
measures. It is the tangent space at the identity of the group of
automorphisms \cite[pp.~677--678]{Hochschild1942}.
Sections~\ref{sec:basic-theory:der-as-lie-algebra}
and~\ref{sec:basic-theory:exponential} tell this part of the story.

The third group studied rings. For them the answer depends on the algebra.
In a matrix algebra over a field, every derivation is a commutator with a fixed
matrix \cite[Corollary~6.2]{BresarHanselkaKlepVolcic2018}. In the Weyl algebra
this can fail, and it depends on the characteristic of the field
\cite[p.~48]{BenkartLopesOndrus2015}. Kharchenko built a theory around such
differences \cite{Kharchenko1991}. Sections~\ref{sec:basic-theory:inner-outer}
and~\ref{sec:basic-theory:extension} follow this theory.

The chapter asks one question: what does the set of all derivations tell us
about the algebra? We go step by step. First comes the definition. Then we
collect all derivations into one object and ask which come from a fixed
element. Next we link derivations to automorphisms. After that we study the
elements that a derivation sends to zero, and we ask when a derivation extends
to a larger algebra.

The last sections look back. We name the close relatives of a derivation and
list the questions that the book asks about each one. The chapter ends with examples from the sources, on
matrices and polynomials.
Each field where these notions apply is reported in the chapter that owns its
notions.

\hbfigure{02-basic-theory/exponential-flow}{The derivation $D = x\,\partial/\partial y - y\,\partial/\partial x$ viewed as a vector field on the plane. On smooth functions, $\exp(tD)$ is the rotation by the angle $t$. On polynomials the series $\exp(tD)$ does not stop, so it is not defined there. Section~\ref{sec:basic-theory:exponential} explains why.}{fig:basic:exp}
